% !PRINTED 713
\weberpage{713}
Beides ergiebt sich leicht dadurch, dass für $n = -1$ im ersten
Falle $\gamma = \tfrac12\mu$, im zweiten $\alpha=1$, $\beta=0$ ist,
also beide Male $\chi_1(-1)=1$ ist, und dass jede der Formeln (6),
(7) genau $\tfrac12\mu$ verschiedene Charaktere darstellt.

Worauf es nun wesentlich ankommt, ist Folgendes:

\begin{quote}
\gesperrt{Ist $p$ irgend eine von $q$ verschiedene Primzahl, so
erhält man in der Form $\chi_1(p)$, wenn $\chi_1$ die Gruppe $C_1$
durchläuft, jede $f_1^{\text{te}}$ Einheitswurzel, und jede gleich
oft, also $e_1$ mal.}
\end{quote}

Dass alle $\chi_1(p)$ Einheitswurzeln vom Grade $f_1$ sind, ist
klar, denn es ist immer (vgl. \S.\,177) $p^{f_1}\equiv\pm1\pmod{m}$,
und daher
\[ [\chi_1(p)]^{f_1} = \chi_1(\pm1) = 1. \]
Wenn aber unter den $\chi_1(p)$ eine primitive $f_1^{\text{te}}$
Einheitswurzel vorkommt, so kommen alle anderen $f_1^{\text{ten}}$
Einheitswurzeln auch darunter vor, und gleich oft, nämlich so oft
wie die Einheitswurzel 1.

Dies folgt unmittelbar aus der Gruppennatur der $\chi_1$. Es ist
also nun noch zu zeigen, dass unter den $\chi_1(p)$ \gesperrt{eine primitive}
$f_1^{\text{te}}$ Einheitswurzel vorkommt.

Diese Möglichkeit ergiebt sich aber aus der Bestimmung von $\chi_1$
sehr einfach. Ist zunächst $\Theta$ in (6) bei \gesperrt{ungeradem $m$ eine
primitive Einheitswurzel vom Grade} $\tfrac12\mu$, und ist $\gamma$
der Index von $p$, so ist eine Potenz von $\Theta^{\gamma}$,
$\Theta^{\gamma\lambda}$ dann und nur dann $=1$, wenn
\[ 2\gamma\lambda \equiv 0 \pmod{\mu}. \]
Nun ist $\mu=ef$ und $e$ der grösste gemeinschaftliche Theiler von
$\gamma$ und $\mu$; und daher wird diese Bedingung
$2\lambda\equiv0\pmod{f}$. Es ist aber nach \S.\,177, 3. $f_1=f$
oder $=\tfrac12 f$, je nachdem $f$ ungerade oder gerade ist, und
folglich muss $\lambda\equiv0\pmod{f_1}$ sein. Also ist
$\Theta^{\gamma}$ primitive $f_1^{\text{te}}$ Einheitswurzel.

Ist aber $m$ \gesperrt{gerade}, so nehmen wir in (7) gleichfalls für $\Theta$
eine primitive $\tfrac12\mu^{\text{te}}$ Einheitswurzel und setzen
\[ p \equiv (-1)^{\alpha}5^{\beta} \pmod{m}, \]
so dass $f$ die kleinste positive Zahl ist, die den Congruenzen
\[ \alpha f \equiv 0 \pmod{2}, \quad \beta f \equiv 0 \pmod{\tfrac12\mu} \]
genügt. Es ist nun $\Theta^{\beta\lambda}$ nur dann $=1$, wenn
$\beta\lambda\equiv0\pmod{\tfrac12\mu}$ ist. Wenn nicht $\lambda=1$,
also $\beta=0$ ist, so ist der kleinste Werth, den $\lambda$ haben
kann, gerade, und daher ist dieser kleinste Werth von $\lambda=f$.
